% Source: https://artofproblemsolving.com/community/c3h1915769_proportional_proportions
% Stable identity: aops-topic-1915769-q01
Write the conditions with positive constants $k_1,k_2,k_3$:
\[
x=\dfrac{k_1y}{z},\qquad z=k_2xy,\qquad w=\dfrac{k_3}{xyz}.
\]
Substituting the second into the first gives $x=\dfrac{k_1y}{k_2xy}=\dfrac{k_1}{k_2x}$, so $x^2=\dfrac{k_1}{k_2}$: the quantity $x$ is constant. Then the third condition gives
\[
yz=\dfrac{k_3}{x}\cdot w^{-1},
\]
so $yz$ is proportional to $w^{-1}$. Since $x$ is constant, $x^{\beta}$ is constant for every integer $\beta$, so $yz$ is proportional to $w^{-1}x^{\beta}$ for \emph{every} integer $\beta$. Such variation really occurs, for example $x=1$, $y=z=t$, $w=\dfrac{1}{t^2}$ for $t>0$. Hence $\alpha+\beta=\beta-1$ can be as large as we like, and the answer is $\boxed{\text{(E)}}$.
